\chapter{Formes différentielles et théorème de Stokes}\label{ch:b3:forms}

Un théorème d'analyse a, en deux siècles, absorbé tous
les autres de son espèce~: le théorème fondamental de l'analyse,
Green--Riemann (prouvé dans le volume d'Année~2), le théorème
de la divergence de Gauss, le théorème du rotationnel de
Kelvin--Stokes --- chacun dit que l'intégrale d'une certaine
dérivée sur une région égale l'intégrale de l'objet original
sur le \omterm{def:b3:forms:boundary}{bord}. Le langage des \emph{\omterm{def:b3:forms:diffform}{formes différentielles}}
en fait un seul énoncé, $\int_M\dd\omega = \int_{\partial
M}\omega$, et rend cet énoncé prouvable d'un seul coup. Ce
chapitre construit le langage honnêtement --- algèbre
multilinéaire alternée, la \omterm{def:b3:forms:d}{dérivée extérieure}, les \omterm{def:b3:forms:pullback}{images
réciproques}, l'\omterm{def:b3:forms:orientation}{orientation}, l'intégration sur les sous-variétés
du~\cref{ch:b3:submanifolds} --- prouve le théorème de Stokes,
et encaisse les premiers chèques~: les théorèmes d'intégrale
classiques, le \emph{\omterm{def:b3:forms:winding}{nombre d'enroulement}} qui faisait tourner
en secret le~\cref{ch:b3:residues}, et, dans le problème de
week-end, le théorème du point fixe de Brouwer. Tout au long,
\emph{lisse} signifie $\mathcal C^\infty$~; toute application
et toute forme est lisse sauf mention contraire. Cela ne coûte
aucune généralité qui vaille à ce niveau et libère les mains.

\section{Algèbre multilinéaire alternée}

\begin{definition}\label{def:b3:forms:alternating}
Soit $E$ un espace vectoriel réel de dimension $n$. Une
\emph{forme $k$-linéaire alternée}\index{forme alternée} sur
$E$ est une application $\alpha\colon E^k \to \R$, linéaire
en chaque variable, avec $\alpha(v_1, \dots, v_k) = 0$ dès
que deux arguments sont égaux. Leur espace s'écrit $\Lambda^k
E^*$~; par convention $\Lambda^0E^* = \R$. L'alternance force
l'antisymétrie~: échanger deux arguments change le signe
(développer $\alpha(\dots, v + w, \dots, v + w, \dots) =
0$), et plus généralement $\alpha(v_{\sigma(1)}, \dots,
v_{\sigma(k)}) = \varepsilon(\sigma)\,\alpha(v_1, \dots,
v_k)$ pour toute permutation $\sigma$.
\end{definition}

\begin{example}
Sur $E = \R^n$~: $\Lambda^1E^* = E^*$ est l'\omterm{def:b3:banach:operator}{espace dual}~; le
déterminant dans la \omterm{def:b3:topology:basis}{base} canonique est une forme $n$-linéaire
alternée, et la~\cref{prop:b3:forms:basis} montrera qu'il engendre
$\Lambda^nE^*$ --- la raison profonde pour laquelle le
déterminant est unique à un scalaire près. Pour $k > n$,
$\Lambda^kE^* = \{0\}$~: $k$ vecteurs sont liés, et
développer l'un le long des autres tue $\alpha$ par
alternance.
\end{example}

\begin{definition}\label{def:b3:forms:wedge}
Pour $\ell_1, \dots, \ell_k \in E^*$, leur \emph{produit
extérieur}\index{produit extérieur} est la forme $k$-linéaire
alternée
\[
(\ell_1 \wedge \dots \wedge \ell_k)(v_1, \dots, v_k)
= \det\bigl(\ell_i(v_j)\bigr)_{1 \leq i, j \leq k} .
\]
Alternance et multilinéarité sont celles du déterminant en
ses colonnes.
\end{definition}

\begin{proposition}[Base de $\Lambda^k$]\label{prop:b3:forms:basis}
Soit $(e_1, \dots, e_n)$ une \omterm{def:b3:topology:basis}{base} de $E$ de \omterm{def:b3:topology:basis}{base} duale
$(e_1^*, \dots, e_n^*)$. Les formes
\[
e_I^* = e_{i_1}^* \wedge \dots \wedge e_{i_k}^*,
\qquad I = \{i_1 < \dots < i_k\} \subseteq \{1, \dots, n\},
\]
forment une \omterm{def:b3:topology:basis}{base} de $\Lambda^kE^*$~; d'où $\dim\Lambda^kE^* =
\binom nk$. Explicitement, $\alpha =
\sum_{\abs I = k}\alpha(e_{i_1}, \dots, e_{i_k})\,e_I^*$.
\end{proposition}

\begin{proof}
\emph{Engendrement.} Soit $\alpha \in \Lambda^kE^*$ et $\beta
= \sum_I\alpha(e_I)\,e_I^*$, où $\alpha(e_I)$ abrège
$\alpha(e_{i_1}, \dots, e_{i_k})$. Les deux côtés sont
$k$-linéaires et alternés, donc ils coïncident dès qu'ils
coïncident sur tous les $k$-uplets $(e_{j_1}, \dots,
e_{j_k})$ avec $j_1 < \dots < j_k$ (la multilinéarité réduit
aux uplets de vecteurs de \omterm{def:b3:topology:basis}{base}, l'alternance aux uplets
strictement croissants). Et $e_I^*(e_{j_1}, \dots, e_{j_k})
= \det(e_{i_r}^*(e_{j_s})) = \delta_{IJ}$~: pour $I = J$ la
matrice est l'identité~; pour $I \neq J$ une ligne est nulle.
Donc $\beta(e_J) = \alpha(e_J)$ pour tout $J$~: $\beta =
\alpha$. \emph{Liberté.} Si $\sum_I c_Ie_I^* = 0$, évaluer
sur $(e_{j_1}, \dots, e_{j_k})$ donne $c_J = 0$.
\end{proof}

\begin{definition}\label{def:b3:forms:wedgegeneral}
Le \omterm{def:b3:forms:wedge}{produit extérieur} s'étend en une application bilinéaire
$\Lambda^kE^* \times \Lambda^\ell E^* \to
\Lambda^{k+\ell}E^*$, déterminée par la bilinéarité et
$(e_I^*) \wedge (e_J^*) = e_I^* \wedge e_J^*$ (concaténer
et réordonner~; le produit est $0$ si $I \cap J \neq
\varnothing$). Elle est associative, et
\emph{anticommutative graduée}~:
\[
\beta \wedge \alpha = (-1)^{k\ell}\,\alpha \wedge \beta
\qquad (\alpha \in \Lambda^k, \ \beta \in \Lambda^\ell).
\]
\end{definition}

\begin{proof}[\omnameProof]
Les deux propriétés se vérifient sur les éléments de \omterm{def:b3:topology:basis}{base} et
s'étendent par bilinéarité. Associativité~: les deux
parenthésages de $e_I^* \wedge e_J^* \wedge e_K^*$ égalent
le wedge de la famille concaténée de $1$-formes, par la
formule du déterminant de la~\cref{def:b3:forms:wedge}
(développement de Laplace par blocs). La règle des signes~:
faire passer chacun des $\ell$ facteurs de $\beta$ devant les
$k$ facteurs de $\alpha$ coûte un signe par transposition
adjacente (échange de deux lignes du déterminant), d'où
$(-1)^{k\ell}$ au total.
\end{proof}

\begin{proposition}[Image réciproque, cas linéaire]
\label{prop:b3:forms:linearpullback}
Une application linéaire $u\colon E \to F$ induit, pour
chaque $k$, l'application linéaire $u^*\colon \Lambda^kF^*
\to \Lambda^kE^*$, $(u^*\alpha)(v_1, \dots, v_k) =
\alpha(u(v_1), \dots, u(v_k))$. Elle vérifie $u^*(\alpha
\wedge \beta) = u^*\alpha \wedge u^*\beta$ et $(u \circ w)^*
= w^* \circ u^*$. De plus~:
\begin{enumerate}
  \item Si $\dim E = n$ et $u\colon E \to E$, alors sur la
        droite $\Lambda^nE^*$~: $u^*\alpha = (\det u)\,\alpha$.
  \item Si $\operatorname{rk}u < k$, alors $u^* = 0$ sur
        $\Lambda^kF^*$.
\end{enumerate}
\end{proposition}

\begin{proof}
Les identités fonctorielles sont immédiates d'après les
définitions (pour la règle du produit, vérifier sur les wedges
de $1$-formes avec la formule du déterminant ---
$\det(\ell_i(uv_j)) = \det((u^*\ell_i)(v_j))$ --- puis
étendre bilinéairement).
(1) $u^*$ envoie la droite $\Lambda^nE^*$
(la~\cref{prop:b3:forms:basis}) dans elle-même, donc $u^*\alpha
= c\,\alpha$ avec $c$ indépendant de $\alpha \neq 0$~; tester
sur $\alpha = e_1^* \wedge \dots \wedge e_n^*$ et $(v_j) =
(e_j)$ donne $c = \det(e_i^*(ue_j)) = \det u$.
(2) Pour $v_1, \dots, v_k \in E$, les vecteurs $u(v_1),
\dots, u(v_k)$ vivent dans l'image de $u$, de dimension $<
k$~: ils sont linéairement dépendants, et une \omterm{def:b3:forms:alternating}{forme alternée}
s'annule sur une famille liée (développer le vecteur lié le
long des autres).
\end{proof}

\section{Formes différentielles et dérivée extérieure}

\begin{definition}\label{def:b3:forms:diffform}
Soit $U \subseteq \R^n$ \omterm{def:b3:topology:topology}{ouvert}. Une \emph{forme
différentielle de degré $k$}\index{forme différentielle} sur
$U$ est une application lisse $\omega\colon U \to
\Lambda^k(\R^n)^*$~; dans la \omterm{def:b3:topology:basis}{base}
de la~\cref{prop:b3:forms:basis} (en écrivant $\dd x_i$ pour
$e_i^*$),
\[
\omega = \sum_{\abs I = k} a_I\,\dd x_I,
\qquad
\dd x_I = \dd x_{i_1} \wedge \dots \wedge \dd x_{i_k},
\]
avec des coefficients lisses $a_I \in \mathcal C^\infty(U)$.
Leur espace est $\Omega^k(U)$~; $\Omega^0(U) = \mathcal
C^\infty(U)$. Une $0$-forme est une fonction~; une $1$-forme
est un champ de formes linéaires (p.ex.\ la différentielle
$\dd f$ d'une fonction)~; une $n$-forme est $a\,\dd
x_1\wedge\dots\wedge\dd x_n$, l'intégrande naturelle
du~\cref{ch:b3:product}.
\end{definition}

\begin{definition}[Dérivée extérieure]\label{def:b3:forms:d}
La \emph{dérivée extérieure}\index{dérivée extérieure} est
l'application linéaire $\dd\colon \Omega^k(U) \to
\Omega^{k+1}(U)$ définie par
\[
\dd\Bigl(\sum_I a_I\,\dd x_I\Bigr)
= \sum_I \dd a_I \wedge \dd x_I
= \sum_I\sum_{j=1}^n \frac{\partial a_I}{\partial
x_j}\,\dd x_j \wedge \dd x_I .
\]
Sur les $0$-formes, c'est la différentielle usuelle.
\end{definition}

\begin{theorem}\label{thm:b3:forms:dsquare}
(a) $\dd(\omega \wedge \eta) = \dd\omega \wedge \eta +
(-1)^k\,\omega \wedge \dd\eta$ pour $\omega \in \Omega^k$
(la \emph{règle de Leibniz graduée}).
(b) $\dd \circ \dd = 0$.\index{d au carré nul@$\dd^2 = 0$}
\end{theorem}

\begin{proof}
(a) Par bilinéarité il suffit de traiter $\omega = a\,\dd
x_I$, $\eta = b\,\dd x_J$. Alors $\omega \wedge \eta =
ab\,\dd x_I \wedge \dd x_J$ et
\[
\dd(\omega\wedge\eta) = (b\,\dd a + a\,\dd b) \wedge \dd x_I
\wedge \dd x_J
= (\dd a \wedge \dd x_I) \wedge (b\,\dd x_J) +
a\,\dd b \wedge \dd x_I \wedge \dd x_J,
\]
et faire passer la $1$-forme $\dd b$ devant les $k$ facteurs
de $\dd x_I$ coûte $(-1)^k$ (la~\cref{def:b3:forms:wedgegeneral})~:
le second terme est $(-1)^k\,\omega \wedge \dd\eta$.
(b) Pour une fonction~: $\dd(\dd a) = \sum_{i,j}
\frac{\partial^2a}{\partial x_j\partial x_i}\,\dd x_j \wedge
\dd x_i$. Le coefficient est symétrique en $(i,j)$ par le
théorème de Schwarz sur les dérivées partielles mixtes
(prouvé dans le volume d'Année~2~; $a$ est $\mathcal
C^\infty$), tandis que $\dd x_j \wedge \dd x_i$ est
antisymétrique~: en appariant les termes $(i,j)$ et $(j,i)$,
tout s'annule. Pour une $\omega = \sum a_I\dd x_I$
générale~: $\dd\dd\omega = \sum\dd(\dd a_I \wedge \dd x_I)
= \sum(\dd\dd a_I\wedge\dd x_I - \dd a_I\wedge\dd(\dd x_I))$
par (a), et les deux termes s'annulent ($\dd(\dd x_I) = 0$
car le coefficient est constant).
\end{proof}

\begin{definition}[Image réciproque]\label{def:b3:forms:pullback}
Soit $\varphi\colon U \to V$ lisse ($U \subseteq \R^m$, $V
\subseteq \R^n$ \omterm{def:b3:topology:topology}{ouverts}). L'\emph{image
réciproque}\index{image réciproque} $\varphi^*\colon
\Omega^k(V) \to \Omega^k(U)$ est définie ponctuellement par
l'image réciproque linéaire le long de la différentielle~:
$(\varphi^*\omega)_x = (D\varphi(x))^*\,\omega_{\varphi(x)}$.
Concrètement, $\varphi^*$ substitue~: $\varphi^*f = f \circ
\varphi$ sur les fonctions, $\varphi^*(\dd y_i) =
\dd\varphi_i = \sum_j\frac{\partial\varphi_i}{\partial
x_j}\dd x_j$, et $\varphi^*(a\,\dd y_{i_1}\wedge\dots\wedge
\dd y_{i_k}) =
(a\circ\varphi)\,\dd\varphi_{i_1}\wedge\dots\wedge
\dd\varphi_{i_k}$.
\end{definition}

\begin{theorem}\label{thm:b3:forms:pullback}
(a) $\varphi^*(\omega \wedge \eta) = \varphi^*\omega \wedge
\varphi^*\eta$ et $(\psi \circ \varphi)^* = \varphi^* \circ
\psi^*$.
(b) $\varphi^*(\dd\omega) = \dd(\varphi^*\omega)$~: la
\omterm{def:b3:forms:d}{dérivée extérieure} commute avec toute substitution lisse ---
l'identité qui en fait \emph{la} dérivée de la théorie.
(c) Si $\varphi\colon U \to V$ est lisse entre \omterm{def:b3:topology:topology}{ouverts} de
$\R^n$ et $\omega = a\,\dd y_1 \wedge \dots \wedge \dd
y_n$, alors $\varphi^*\omega = (a \circ
\varphi)\,\det\bigl(D\varphi\bigr)\,\dd x_1 \wedge \dots
\wedge \dd x_n$.
\end{theorem}

\begin{proof}
(a) Énoncés ponctuels sur les \omterm{def:b3:forms:pullback}{images réciproques} linéaires
(la~\cref{prop:b3:forms:linearpullback}), plus la règle de
chaîne $D(\psi\circ\varphi)(x) =
D\psi(\varphi(x))\,D\varphi(x)$.
(b) Pour une $0$-forme $f$~: $\varphi^*(\dd f) = \dd f \circ
D\varphi = \dd(f \circ \varphi)$ est la règle de chaîne.
Pour $\omega = a\,\dd y_I$~: en utilisant (a),
$\varphi^*\omega =
(a\circ\varphi)\,\dd\varphi_{i_1}\wedge\dots\wedge
\dd\varphi_{i_k}$, donc par la règle de Leibniz
(le~\cref{thm:b3:forms:dsquare}(a)) et $\dd\dd\varphi_{i_r} =
0$~:
\[
\dd(\varphi^*\omega) = \dd(a\circ\varphi) \wedge
\dd\varphi_{i_1}\wedge\dots\wedge\dd\varphi_{i_k}
= \varphi^*(\dd a) \wedge \varphi^*(\dd y_I)
= \varphi^*(\dd a \wedge \dd y_I) =
\varphi^*(\dd\omega).
\]
(c) Ponctuellement c'est exactement $u^*\alpha = (\det
u)\alpha$ sur les formes de \omterm{def:b3:galois:extension}{degré} \omterm{def:b3:rings:primemaximal}{maximal}
(la~\cref{prop:b3:forms:linearpullback}(1) avec $u =
D\varphi(x)$).
\end{proof}

\begin{example}[Coordonnées polaires]\label{ex:b3:forms:polar}
Pour $\varphi(r, \theta) = (r\cos\theta, r\sin\theta)$~:
$\varphi^*\dd x = \cos\theta\,\dd r -
r\sin\theta\,\dd\theta$, $\varphi^*\dd y = \sin\theta\,\dd r
+ r\cos\theta\,\dd\theta$, donc
\[
\varphi^*(\dd x \wedge \dd y) =
(\cos\theta\,\dd r - r\sin\theta\,\dd\theta) \wedge
(\sin\theta\,\dd r + r\cos\theta\,\dd\theta)
= r\,\dd r \wedge \dd\theta,
\]
le jacobien de l'\cref{ex:b3:product:polar} apparaissant par
pure algèbre --- sans théorie de la \omterm{def:b3:measure:measure}{mesure}.
l'\cref{exo:b3:forms:9} transforme cette remarque en énoncé~:
pour les intégrales orientées, la formule de changement de
variables \emph{est} la formule d'\omterm{def:b3:forms:pullback}{image réciproque}.
\end{example}

\section{Formes fermées et exactes~; le lemme de Poincaré}

\begin{definition}\label{def:b3:forms:closedexact}
$\omega \in \Omega^k(U)$ est \emph{fermée}\index{forme
fermée} si $\dd\omega = 0$, \emph{exacte}\index{forme
exacte} si $\omega = \dd\eta$ pour une $\eta \in
\Omega^{k-1}(U)$ (une \emph{primitive} de $\omega$). Exacte
$\Rightarrow$ fermée par $\dd^2 = 0$~; la réciproque est une
question sur la \emph{forme} de $U$.
\end{definition}

\begin{example}[La forme angulaire]\label{ex:b3:forms:angular}
Sur $U = \R^2 \setminus \{0\}$,
\[
\omega_\theta = \frac{x\,\dd y - y\,\dd x}{x^2 + y^2}
\]
est \omterm{def:b3:forms:closedexact}{fermée} (calcul direct~: l'\cref{exo:b3:forms:4}) mais
pas \omterm{def:b3:forms:closedexact}{exacte}~: son intégrale le long du cercle unité vaut $2\pi
\neq 0$, tandis que les intégrales de \omterm{def:b3:forms:closedexact}{formes exactes} le long
de courbes \omterm{def:b3:forms:closedexact}{fermées} s'annulent
(la~\cref{prop:b3:forms:exactloop}). Localement, $\omega_\theta
= \dd\theta$ pour toute détermination lisse $\theta$ de
l'angle polaire --- d'où le nom et l'obstruction~: aucune
telle détermination n'existe sur tout $U$. Cette seule forme
fait tourner le \omterm{def:b3:forms:winding}{nombre d'enroulement}
(la~\cref{sec:b3:forms:winding}) et, à travers lui, le théorème
des \omterm{def:b3:residues:singularities}{résidus} du~\cref{ch:b3:residues}.
\end{example}

\begin{theorem}[Lemme de Poincaré]\label{thm:b3:forms:poincare}
Soit $U \subseteq \R^n$ \omterm{def:b3:topology:topology}{ouvert} et \emph{étoilé} par rapport
à $0$. Toute $k$-forme \omterm{def:b3:forms:closedexact}{fermée} sur $U$ ($k \geq 1$) est
\omterm{def:b3:forms:closedexact}{exacte}.\index{lemme de Poincaré}
\end{theorem}

\begin{proof}
On construit un \emph{opérateur d'homotopie} linéaire
$h\colon \Omega^k(U) \to \Omega^{k-1}(U)$ avec
\begin{equation}\label{eq:b3:forms:homotopy}
\dd(h\omega) + h(\dd\omega) = \omega
\qquad (k \geq 1);
\end{equation}
si $\dd\omega = 0$, alors $\omega = \dd(h\omega)$ et on a
fini. Pour $\omega = \sum_I a_I\,\dd x_I$ poser
\[
h\omega = \sum_{\abs I = k}\ \sum_{r=1}^{k}(-1)^{r-1}
\Bigl(\int_0^1 t^{k-1}a_I(tx)\,\dd t\Bigr)\,
x_{i_r}\,\dd x_{i_1}\wedge\dots\wedge
\widehat{\dd x_{i_r}}\wedge\dots\wedge\dd x_{i_k}
\]
(le chapeau supprime un facteur~; les intégrales sont lisses
en $x$ par dérivation sous le signe intégral,
le~\cref{thm:b3:lebesgue:paramdiff}, toutes les dérivées étant
dominées sur les compacts). Vérifier
\eqref{eq:b3:forms:homotopy} est un calcul fait une fois
dans une vie, donc on le fait en entier. Fixer $I$ et prendre
$\omega = a\,\dd x_I$ (linéarité). D'abord,
\[
\dd(h\omega) = k\Bigl(\int_0^1t^{k-1}a(tx)\dd t\Bigr)\dd x_I
+ \sum_{r=1}^k(-1)^{r-1}\sum_{j=1}^n
\Bigl(\int_0^1t^{k}\,\partial_ja(tx)\,\dd t\Bigr)
x_{i_r}\,\dd x_j\wedge\dd x_{I\setminus i_r} :
\]
le premier groupe rassemble les termes où $\dd$ frappe le
facteur $x_{i_r}$ --- le wedge $\dd x_{i_r}\wedge\dd
x_{I\setminus i_r}$ réassemble $\dd x_I$ avec un signe
$(-1)^{r-1}$ qui annule le préfacteur, et les $k$ valeurs de
$r$ donnent le facteur $k$ --- tandis que le second groupe
rassemble les termes où $\dd$ frappe l'intégrale (la règle
de chaîne sort $t\,\partial_ja(tx)$). Ensuite, $\dd\omega =
\sum_j\partial_ja\,\dd x_j\wedge\dd x_I$, et en appliquant
la définition de $h$ en \omterm{def:b3:galois:extension}{degré} $k+1$, l'\omterm{def:b3:holomorphic:index}{indice} $j$ occupant
la première place~:
\[
h(\dd\omega) = \sum_{j=1}^n\Bigl(\int_0^1t^{k}
\partial_ja(tx)\dd t\Bigr)x_j\,\dd x_I
- \sum_{j=1}^n\sum_{r=1}^k(-1)^{r-1}
\Bigl(\int_0^1t^{k}\partial_ja(tx)\dd t\Bigr)
x_{i_r}\,\dd x_j\wedge\dd x_{I\setminus i_r} .
\]
Les doubles sommes s'annulent dans $\dd(h\omega) +
h(\dd\omega)$, qui égale donc
\[
\Bigl(\int_0^1\bigl(k\,t^{k-1}a(tx) +
t^k{\textstyle\sum_j}x_j\,\partial_ja(tx)\bigr)\dd t\Bigr)
\dd x_I
= \Bigl(\int_0^1\frac{\dd}{\dd t}\bigl(t^ka(tx)\bigr)\dd
t\Bigr)\dd x_I
= a(x)\,\dd x_I = \omega,
\]
par le théorème fondamental de l'analyse. L'étoilé est entré
là où il devait~: $tx \in U$ pour $t \in \intcc01$, de sorte
que $a(tx)$ a un sens.
\end{proof}

\begin{remark}
Pour $k = 1$ et $\omega = \sum_j a_j\dd x_j$, la primitive
est $f(x) = \int_0^1\sum_ja_j(tx)\,x_j\,\dd t$ ---
l'intégrale curviligne de $\omega$ le long du segment $[0,
x]$~: le théorème est le «~un champ à jacobienne symétrique
est un gradient~» à plusieurs variables de l'Année~2,
maintenant en tout \omterm{def:b3:galois:extension}{degré}. La forme angulaire
(l'\cref{ex:b3:forms:angular}) montre que l'hypothèse sur $U$
n'est pas décorative~: $\R^2\setminus\{0\}$ n'est pas
étoilé, et là fermeture n'implique pas exactitude. Ce qui
survit sur un \omterm{def:b3:topology:topology}{ouvert} général est mesuré par la
\emph{cohomologie de de Rham} $H^k(U) =
\ker\dd/\operatorname{im}\dd$ --- voir
l'\cref{exo:b3:forms:12} pour le premier calcul non trivial.
\end{remark}

\section{Orientation et intégration sur les sous-variétés}

Intégrer une $k$-forme demande un territoire $k$-dimensionnel
orienté. Rappeler du~\cref{ch:b3:submanifolds}
(le~\cref{thm:b3:submanifolds:characterizations}) qu'une
$k$-sous-variété $M \subseteq \R^n$ est localement l'image
d'une paramétrisation régulière $\gamma\colon V \to M \cap
W$ ($V \subseteq \R^k$ \omterm{def:b3:topology:topology}{ouvert}, $\gamma$ un \omterm{def:b3:topology:continuity}{homéomorphisme} sur
son image avec différentielle injective).

\begin{definition}\label{def:b3:forms:orientation}
Une \emph{orientation}\index{orientation} de $M$ est un
choix, pour chaque $p \in M$, d'une des deux classes
d'orientation de \omterm{def:b3:topology:basis}{bases} de l'\omterm{def:b3:submanifolds:tangent}{espace tangent} $T_pM$, qui est
\emph{localement cohérent}~: autour de chaque point il existe
une paramétrisation $\gamma$ dont le repère de coordonnées
$(\partial_1\gamma, \dots, \partial_k\gamma)$ est
positivement orienté en tout point de son domaine. De telles
paramétrisations sont dites \emph{directes}. $M$ est
\emph{orientable} si une orientation existe~; le ruban de
Möbius montre que cela peut échouer. Toutes les sous-variétés
de ce chapitre sont orientées.
\end{definition}

\begin{definition}[Intégrale d'une forme]\label{def:b3:forms:integral}
Soit $M$ une $k$-sous-variété orientée et $\omega$ une
$k$-forme définie sur un \omterm{def:b3:topology:topology}{voisinage} de $M$, avec
$\operatorname{supp}\omega \cap M$ compact.
(a) Si $\operatorname{supp}\omega \cap M \subseteq
\gamma(V)$ pour une seule paramétrisation directe, poser
\[
\int_M\omega = \int_V\gamma^*\omega
\]
--- le second membre étant l'intégrale de Lebesgue sur $V$
(le~\cref{ch:b3:product}) du coefficient $g$ de
$\gamma^*\omega = g\,\dd u_1\wedge\dots\wedge\dd u_k$, qui
est \omterm{def:b3:topology:continuity}{continue} à support compact.
(b) En général, choisir un nombre fini de paramétrisations
directes $\gamma_i(V_i)$ recouvrant le compact
$\operatorname{supp}\omega \cap M$ et une
\emph{\omterm{lem:b3:forms:partition}{partition de l'unité}} subordonnée $(\chi_i)$
(le~\cref{lem:b3:forms:partition}), et poser $\int_M\omega =
\sum_i\int_M\chi_i\,\omega$, chaque terme calculé par (a).
Pour une courbe ($k = 1$) paramétrée par
$\gamma\colon\intcc ab\to\R^n$ on écrit
$\int_\gamma\omega = \int_a^b\gamma^*\omega$, sans exiger
l'injectivité.
\end{definition}

\begin{lemma}[Cohérence]\label{lem:b3:forms:consistency}
La définition (a) ne dépend pas de la paramétrisation
directe, et la définition (b) ne dépend ni du recouvrement ni
de la \omterm{lem:b3:forms:partition}{partition de l'unité}. De plus, si $\Phi$ est un
difféomorphisme de \omterm{def:b3:topology:topology}{voisinages} de deux sous-variétés orientées
avec $\Phi(M) = M'$, transportant un repère direct de $M$
vers un repère direct de $M'$ en un point de chaque composante
de $M$, alors $\int_{M'}\omega = \int_M\Phi^*\omega$.
\end{lemma}

\begin{proof}
(a) Soient $\gamma\colon V \to M$, $\delta\colon V' \to M$
des paramétrisations directes dont les images contiennent
$\operatorname{supp}\omega\cap M$. La transition $\tau =
\delta^{-1}\circ\gamma$ est un difféomorphisme entre les
\omterm{def:b3:topology:topology}{ouverts} pertinents de $V, V'$ (lissité des transitions~:
le~\cref{thm:b3:submanifolds:characterizations}, via la
description en graphe local), et $\gamma = \delta\circ\tau$
là, donc $\gamma^*\omega = \tau^*(\delta^*\omega)$
(le~\cref{thm:b3:forms:pullback}(a)). Écrire $\delta^*\omega =
g\,\dd u_1\wedge\dots\wedge\dd u_k$~; alors
(le~\cref{thm:b3:forms:pullback}(c)) $\tau^*(\delta^*\omega) =
(g\circ\tau)\,\det(D\tau)\,\dd u_1\wedge\dots\wedge\dd u_k$.
Les deux repères étant directs, $D\tau$ envoie une \omterm{def:b3:topology:basis}{base}
positive sur une \omterm{def:b3:topology:basis}{base} positive~: $\det D\tau > 0$, donc
$\det D\tau = \abs{\det D\tau}$ et le théorème de changement
de variables (le~\cref{thm:b3:product:changeofvar}) donne
$\int(g\circ\tau)\abs{\det D\tau} = \int g$~: les deux
intégrales coïncident. \emph{Voici toute la raison d'être de
l'\omterm{def:b3:forms:orientation}{orientation}}~: sans le contrôle du signe, le jacobien et sa
valeur absolue diffèrent et l'intégrale est mal définie.
(b) Si $(\chi_i)$ et $(\tilde\chi_j)$ sont deux partitions
admissibles (recouvrements inclus), alors par (a) et
l'additivité finie, $\sum_i\int\chi_i\omega =
\sum_{i,j}\int\chi_i\tilde\chi_j\omega =
\sum_j\int\tilde\chi_j\omega$, chaque double terme calculable
dans l'une ou l'autre carte. Le dernier énoncé~: si $\gamma$
parcourt les paramétrisations directes de $M$, alors
$\Phi\circ\gamma$ parcourt les paramétrisations directes de
$M'$ (la comparaison d'\omterm{def:b3:forms:orientation}{orientation} est localement constante,
et fixée en un point par composante), et $(\Phi\circ\gamma)^*\omega
= \gamma^*(\Phi^*\omega)$.
\end{proof}

\begin{lemma}[Partitions de l'unité, cas compact]
\label{lem:b3:forms:partition}
Soit $K \subseteq \R^n$ compact et $W_1, \dots, W_m$ des
\omterm{def:b3:topology:topology}{ouverts} recouvrant $K$. Il existe $\chi_1, \dots, \chi_m \in
\mathcal C^\infty(\R^n)$ avec $0 \leq \chi_i \leq 1$,
$\operatorname{supp}\chi_i \subseteq W_i$ compact, et
$\sum\chi_i = 1$ sur un \omterm{def:b3:topology:topology}{voisinage} de
$K$.\index{partition de l'unité}
\end{lemma}

\begin{proof}
Chaque $x \in K$ vit dans un certain $W_{i(x)}$ avec une
boule \omterm{def:b3:forms:closedexact}{fermée} $\bar B(x, 2r_x) \subseteq W_{i(x)}$~; la
compacité extrait $x_1, \dots, x_N$ avec les boules $B(x_s,
r_{x_s})$ recouvrant $K$. Pour chaque $s$ prendre une bosse
$\theta_s \in \mathcal C^\infty$, $0 \leq \theta_s \leq 1$,
$\theta_s = 1$ sur $\bar B(x_s, r_{x_s})$,
$\operatorname{supp}\theta_s \subseteq B(x_s, 2r_{x_s})$
(régulariser l'indicatrice de la boule de rayon
$\frac32r_{x_s}$, le~\cref{thm:b3:lp:regularization}). Assigner
chaque $s$ à un \omterm{def:b3:holomorphic:index}{indice} $i(s)$ avec $B(x_s, 2r_{x_s})
\subseteq W_{i(s)}$ et poser $\Theta_i = \sum_{i(s) =
i}\theta_s$. Sur l'\omterm{def:b3:topology:topology}{ouvert} $\Omega_0 = \{\sum_j\Theta_j >
\tfrac12\} \supseteq K$, les fonctions
$\Theta_i/\sum_j\Theta_j$ font le travail mais ne sont
définies que là~; pour globaliser, soit $\rho \in \mathcal
C^\infty(\R^n)$ vérifiant $\rho = 0$ là où $\sum_j\Theta_j
\geq 1$ et $\rho > 0$ là où $\sum_j\Theta_j \leq \tfrac12$
(régulariser une coupure convenable de $1 - \sum_j\Theta_j$),
et poser
\[
\chi_i = \frac{\Theta_i}{\rho + \sum_j\Theta_j} .
\]
Le dénominateur est partout $> 0$ et égale $\sum_j\Theta_j$
sur $\{\sum_j\Theta_j \geq 1\}$, un \omterm{def:b3:topology:topology}{voisinage} \omterm{def:b3:topology:topology}{ouvert} de $K$
(chaque point de $K$ vit dans une boule $B(x_s, r_{x_s})$ où
$\theta_s = 1$)~; là $\sum_i\chi_i = 1$. Supports et bornes
sont clairs.
\end{proof}

\section{Le théorème de Stokes}

\begin{definition}\label{def:b3:forms:boundary}
Une \emph{$k$-sous-variété à bord}\index{sous-variété à
bord} $M \subseteq \R^n$ est un ensemble recouvert par des
paramétrisations régulières de deux sortes~: des
\emph{cartes intérieures} $\gamma\colon V \to M \cap W$
avec $V \subseteq \R^k$ \omterm{def:b3:topology:topology}{ouvert}, et des \emph{cartes de bord}
$\gamma\colon V \cap H^k \to M \cap W$, où $H^k = \{u \in
\R^k : u_k \geq 0\}$ et $\gamma$ s'étend de façon lisse et
régulière à l'\omterm{def:b3:topology:topology}{ouvert} $V$. Le \emph{bord} $\partial M$ est
l'ensemble des points atteints en $u_k = 0$~; c'est une
$(k-1)$-sous-variété sans bord, paramétrée par les
applications $u' \mapsto \gamma(u', 0)$. Une \omterm{def:b3:forms:orientation}{orientation} de
$M$ en induit une sur $\partial M$ par la règle de la
\emph{normale sortante d'abord}~: en $p \in \partial M$, une
\omterm{def:b3:topology:basis}{base} $(w_1, \dots, w_{k-1})$ de $T_p\partial M$ est positive
ssi $(\nu, w_1, \dots, w_{k-1})$ est une \omterm{def:b3:topology:basis}{base} positive de
$T_pM$, où $\nu \in T_pM \setminus T_p\partial M$ pointe
hors de $M$ (dans une carte de bord~: $\nu =
-\partial_k\gamma$, à des composantes tangentielles près ---
la classe d'\omterm{def:b3:forms:orientation}{orientation} ne les voit pas).
\end{definition}

\begin{omfigure}
\begin{tikzpicture}[scale=1.05]
  \fill[omDef!10] (-2,0) arc (180:0:2) -- cycle;
  \draw[thick, omDef] (-2,0) arc (180:0:2);
  \draw[thick, omThm] (-2,0) -- (2,0);
  \node[omDef] at (0.1,1.35) {$M$};
  \node[omThm, below] at (1.55,-0.05) {$\partial M$};
  \draw[->, thick, omExo] (1.415,1.415) -- (2.05,2.05);
  \node[omExo, right] at (2.02,2.12) {$\nu$};
  \draw[->, thick] (1.29,1.53) arc (50:75:2);
  \draw[->, thick, omExo] (-0.6,0) -- (-0.6,-0.65);
  \node[omExo, below] at (-0.6,-0.65) {$\nu$};
  \draw[->, thick] (-0.35,0) -- (0.55,0);
  \node[below] at (0.1,-0.14) {induite};
\end{tikzpicture}

{\small La règle de la normale sortante d'abord~: en chaque
point du \omterm{def:b3:forms:boundary}{bord}, placer le vecteur sortant $\nu$ en premier~;
les \omterm{def:b3:topology:basis}{bases} qui le complètent en un repère positif de $M$
orientent $\partial M$. Pour un domaine plan muni de
l'\omterm{def:b3:forms:orientation}{orientation} standard, c'est la règle du sens trigonométrique
de Green--Riemann.}
\end{omfigure}

\begin{lemma}[Stokes sur le demi-espace]
\label{lem:b3:forms:halfspace}
Soit $\eta$ une $(k-1)$-forme lisse sur $\R^k$ à support
compact. Alors
\[
\int_{H^k}\dd\eta = \int_{\partial H^k}\eta,
\]
où $H^k$ porte l'\omterm{def:b3:forms:orientation}{orientation} standard de $\R^k$ et $\partial
H^k = \{u_k = 0\} \cong \R^{k-1}$ l'\omterm{def:b3:forms:orientation}{orientation} induite, qui
est $(-1)^k$ fois l'\omterm{def:b3:forms:orientation}{orientation} standard de $\R^{k-1}$.
\end{lemma}

\begin{proof}
D'abord la comptabilité d'\omterm{def:b3:forms:orientation}{orientation}~: la normale sortante
le long de $\partial H^k$ est $-e_k$, et
\[
\det(-e_k, e_1, \dots, e_{k-1}) =
-\det(e_k, e_1, \dots, e_{k-1}) = -(-1)^{k-1}\det(e_1,
\dots, e_k) = (-1)^k :
\]
le repère $(e_1, \dots, e_{k-1})$ de $\partial H^k$ est
positif pour l'\omterm{def:b3:forms:orientation}{orientation} induite exactement lorsque $k$ est
pair, d'où la comparaison annoncée. Par linéarité prendre
$\eta = f\,\dd u_1\wedge\dots\wedge\widehat{\dd
u_i}\wedge\dots\wedge\dd u_k$, $f \in \mathcal
C^\infty_c(\R^k)$~; alors $\dd\eta =
(-1)^{i-1}\,\partial_if\,\dd u_1\wedge\dots\wedge\dd u_k$
(placer $\dd u_i$ dans son créneau coûte $i - 1$ échanges).
Deux cas, tous deux par Tonelli--Fubini
(le~\cref{thm:b3:product:fubini}) et le théorème fondamental
à une variable.

\emph{Cas $i < k$.} En intégrant d'abord en $u_i$ sur $\R$~:
$\int_\R\partial_if\,\dd u_i = 0$ (support compact), donc
$\int_{H^k}\dd\eta = 0$. Et la restriction de $\eta$ à
$\{u_k = 0\}$ contient le facteur $\dd u_k$, qui se restreint
à $0$ ($u_k$ y est constant)~: $\int_{\partial H^k}\eta = 0$
aussi.

\emph{Cas $i = k$.} En intégrant d'abord en $u_k$ sur
$\intco0\infty$~:
\begin{align*}
\int_{H^k}\dd\eta &= (-1)^{k-1}\int_{\R^{k-1}}
\Bigl(\int_0^\infty\partial_kf\,\dd u_k\Bigr)\dd u' \\
&= (-1)^{k-1}\int_{\R^{k-1}}\bigl(0 - f(u', 0)\bigr)\dd u'
= (-1)^k\int_{\R^{k-1}}f(u',0)\,\dd u' .
\end{align*}
Du côté du \omterm{def:b3:forms:boundary}{bord}, $\eta$ se restreint à $f(u', 0)\,\dd
u_1\wedge\dots\wedge\dd u_{k-1}$, et l'\omterm{def:b3:forms:orientation}{orientation} induite
étant $(-1)^k$ fois l'\omterm{def:b3:forms:orientation}{orientation} standard,
$\int_{\partial H^k}\eta =
(-1)^k\int_{\R^{k-1}}f(u',0)\,\dd u'$. Les deux côtés
coïncident.
\end{proof}

\begin{theorem}[Stokes]\label{thm:b3:forms:stokes}
Soit $M \subseteq \R^n$ une $k$-sous-variété orientée compacte
à \omterm{def:b3:forms:boundary}{bord}, $\partial M$ portant l'\omterm{def:b3:forms:orientation}{orientation} induite, et soit
$\omega$ une $(k-1)$-forme lisse sur un \omterm{def:b3:topology:topology}{voisinage} de $M$.
Alors\index{théorème de Stokes}
\[
\int_M\dd\omega = \int_{\partial M}\omega .
\]
En particulier, si $\partial M = \varnothing$~:
$\int_M\dd\omega = 0$.
\end{theorem}

\begin{proof}
Recouvrir le compact $M$ par un nombre fini d'images de
cartes directes (intérieures ou de \omterm{def:b3:forms:boundary}{bord}), et prendre une
\omterm{lem:b3:forms:partition}{partition de l'unité} $(\chi_i)$ subordonnée aux \omterm{def:b3:topology:topology}{ouverts}
correspondants $W_i$ de $\R^n$ (le~\cref{lem:b3:forms:partition}),
avec $\sum\chi_i = 1$ sur un \omterm{def:b3:topology:topology}{voisinage} de $M$. Sur ce
\omterm{def:b3:topology:topology}{voisinage} $\dd(\sum\chi_i) = 0$, donc
\[
\int_M\dd\omega = \sum_i\int_M\dd(\chi_i\omega),
\qquad
\int_{\partial M}\omega =
\sum_i\int_{\partial M}\chi_i\omega :
\]
les deux côtés sont additifs, et il suffit de prouver le
théorème pour une forme supportée dans l'image d'une seule
carte.

\emph{Carte intérieure.} Si $\operatorname{supp}\omega \cap
M \subseteq \gamma(V)$ avec $V$ \omterm{def:b3:topology:topology}{ouvert} dans $\R^k$~: étendre
$\eta = \gamma^*\omega$ par zéro à $\R^k$ (lisse, support
compact dans $V$) et appliquer le~\cref{lem:b3:forms:halfspace}
avec le support loin de $\partial H^k$ (translater $V$ dans
le demi-espace supérieur \omterm{def:b3:topology:topology}{ouvert} --- ou simplement répéter le
calcul du cas $i < k$ sur tout $\R^k$)~: $\int_M\dd\omega =
\int_{\R^k}\dd\eta = 0$, et $\int_{\partial M}\omega = 0$
puisque $\omega$ s'annule près de $\partial M$.

\emph{Carte de \omterm{def:b3:forms:boundary}{bord}.} Si $\operatorname{supp}\omega\cap M
\subseteq \gamma(V \cap H^k)$~: avec $\eta = \gamma^*\omega$
étendue par zéro, $\gamma^*(\dd\omega) = \dd\eta$
(le~\cref{thm:b3:forms:pullback}(b)), donc par
la~\cref{def:b3:forms:integral}~:
\[
\int_M\dd\omega = \int_{H^k}\dd\eta
\overset{\text{\cref{lem:b3:forms:halfspace}}}{=}
\int_{\partial H^k}\eta .
\]
Il reste à identifier le second membre avec $\int_{\partial
M}\omega$. Le \omterm{def:b3:forms:boundary}{bord} est paramétré par $\beta(u') =
\gamma(u', 0)$, et $\beta^*\omega$ est la restriction de
$\eta$ à $\{u_k = 0\}$ (\omterm{def:b3:forms:pullback}{image réciproque} sous l'inclusion
$u' \mapsto (u', 0)$ composée avec $\gamma$). La comparaison
d'\omterm{def:b3:forms:orientation}{orientation} est le même $(-1)^k$ des deux côtés~: le repère
$(\partial_1\beta, \dots, \partial_{k-1}\beta)$ siège dans
l'\omterm{def:b3:forms:orientation}{orientation} induite de $\partial M$ avec le signe
$\det(-e_k, e_1, \dots, e_{k-1}) = (-1)^k$ calculé dans la
carte (le vecteur sortant se tire en arrière en $-e_k$), qui
est exactement le signe reliant l'\omterm{def:b3:forms:orientation}{orientation} induite de
$\partial H^k$ à l'\omterm{def:b3:forms:orientation}{orientation} standard de $\R^{k-1}$
(le~\cref{lem:b3:forms:halfspace}). Les deux conventions de
signe s'annulent~: $\int_{\partial H^k}\eta =
\int_{\partial M}\omega$.
\end{proof}

\begin{example}[Les théorèmes classiques]
\label{ex:b3:forms:classical}
Soit $D \subseteq \R^2$ un domaine compact à \omterm{def:b3:forms:boundary}{bord} lisse,
orienté standardement. Pour $\omega = P\,\dd x + Q\,\dd y$~:
$\dd\omega = \bigl(\partial_xQ - \partial_yP\bigr)\dd
x\wedge\dd y$, et Stokes se lit
\[
\int_D\Bigl(\frac{\partial Q}{\partial x} -
\frac{\partial P}{\partial y}\Bigr)\dd x\,\dd y
= \oint_{\partial D}P\,\dd x + Q\,\dd y :
\]
\emph{Green--Riemann}, prouvé pour des domaines élémentaires
dans le volume d'Année~2 et maintenant en généralité
naturelle. Dans $\R^3$, Stokes appliqué à la $2$-forme de
flux d'un champ de vecteurs sur un domaine compact donne le
\emph{théorème de la divergence}
$\int_\Omega\operatorname{div}F =
\int_{\partial\Omega}\langle F, \nu\rangle\,\dd S$, et
appliqué à une $1$-forme sur une surface à \omterm{def:b3:forms:boundary}{bord}, le
\emph{théorème de Kelvin--Stokes} classique sur le
rotationnel~; l'\cref{exo:b3:forms:7} explicite les deux
dictionnaires.
\end{example}

\section{Le nombre d'enroulement}\label{sec:b3:forms:winding}

\begin{definition}\label{def:b3:forms:winding}
Soit $\gamma\colon\intcc01\to\R^2\setminus\{a\}$ une courbe
lisse \omterm{def:b3:forms:closedexact}{fermée}. Son \emph{nombre
d'enroulement}\index{nombre d'enroulement} autour de $a$ est
\[
\operatorname{Ind}_\gamma(a) =
\frac1{2\pi}\int_\gamma\omega_\theta^a,
\qquad
\omega_\theta^a = \frac{(x - a_1)\,\dd y - (y - a_2)\,\dd
x}{(x - a_1)^2 + (y - a_2)^2} .
\]
\end{definition}

\begin{proposition}\label{prop:b3:forms:windinginteger}
$\operatorname{Ind}_\gamma(a) \in \Z$~; comme fonction de
$a$ elle est constante sur chaque \omterm{def:b3:topology:components}{composante connexe} de
$\R^2\setminus\gamma(\intcc01)$ et nulle sur la composante
non bornée. Pour $\gamma(t) = a + r(\cos2\pi t, \sin2\pi
t)$~: $\operatorname{Ind}_\gamma(a) = 1$.
\end{proposition}

\begin{proof}
Prendre $a = 0$ et écrire $\gamma = (x, y)$, $\rho =
\norm\gamma > 0$, $\vartheta(t) =
\int_0^t\gamma^*\omega_\theta$, de sorte que $\vartheta' =
\frac{xy' - yx'}{x^2 + y^2}$. En notation complexe soit
$u(t) = \gamma(t)\,\rho(t)^{-1}\eu^{-\iu\vartheta(t)}$~;
alors $\abs u = 1$ et un calcul direct donne
\[
\frac{u'}{u} = \frac{\gamma'}{\gamma} - \frac{\rho'}{\rho}
- \iu\vartheta'
= \frac{\bar\gamma\gamma'}{\abs\gamma^2} -
\frac{\rho'}{\rho} - \iu\,\frac{xy' -
yx'}{\rho^2}
= \frac{(xx' + yy')}{\rho^2} - \frac{\rho'}{\rho} = 0,
\]
puisque $\bar\gamma\gamma' = (xx' + yy') + \iu(xy' - yx')$ et
$\rho\rho' = xx' + yy'$. Donc $u$ est constante~: $\gamma(t)
= \rho(t)\,c\,\eu^{\iu\vartheta(t)}$ avec $\abs c = 1$, et
$\gamma(1) = \gamma(0)$ avec $\rho(1) = \rho(0)$ force
$\eu^{\iu\vartheta(1)} = \eu^{\iu\vartheta(0)} = 1$~:
$\vartheta(1) \in 2\pi\Z$, i.e.\
$\operatorname{Ind}_\gamma(0) \in \Z$. Comme fonction de $a$
sur le complémentaire \omterm{def:b3:topology:topology}{ouvert} de la courbe compacte,
l'intégrale définissante est \omterm{def:b3:topology:continuity}{continue}
(le~\cref{thm:b3:lebesgue:paramcont}, domination sur un
\omterm{def:b3:topology:topology}{voisinage} de chaque $a$)~; une fonction \omterm{def:b3:topology:continuity}{continue} à valeurs
entières est localement constante, donc constante sur les
composantes. Pour $\norm a$ grand l'intégrande est
$O(1/\norm a)$ uniformément en $t$, donc l'\omterm{def:b3:holomorphic:index}{indice} tend vers
$0$ et s'annule sur la composante non bornée. Pour le
cercle~: $\gamma^*\omega_\theta^a = 2\pi\,\dd t$
directement.
\end{proof}

\begin{remark}
Sous $\C \cong \R^2$, $\frac{\dd z}{z} = \frac{x\,\dd x +
y\,\dd y}{x^2 + y^2} + \iu\,\omega_\theta$, donc
$\operatorname{Ind}_\gamma(a) =
\frac1{2\iu\pi}\oint_\gamma\frac{\dd z}{z - a}$~: c'est
l'\omterm{def:b3:holomorphic:index}{indice} du~\cref{ch:b3:residues}, et
la~\cref{prop:b3:forms:windinginteger} reprouve son
intégralité et sa constance locale par des moyens de
variables réelles --- la moitié topologique du théorème des
\omterm{def:b3:residues:singularities}{résidus}, maintenant appuyée sur Stokes.
\end{remark}

\begin{proposition}\label{prop:b3:forms:exactloop}
Si $\omega = \dd f$ est \omterm{def:b3:forms:closedexact}{exacte} sur l'\omterm{def:b3:topology:topology}{ouvert} $U$ et
$\gamma\colon\intcc01\to U$ est une courbe \omterm{def:b3:forms:closedexact}{fermée}, alors
$\int_\gamma\omega = 0$.
\end{proposition}

\begin{proof}
$\int_\gamma\dd f = \int_0^1(f\circ\gamma)'(t)\,\dd t =
f(\gamma(1)) - f(\gamma(0)) = 0$ --- la règle de chaîne
identifie $\gamma^*(\dd f)$ à $(f\circ\gamma)'\,\dd t$.
\end{proof}

\begin{method}\label{met:b3:forms:method}
Calculer avec des formes~: (1) mécaniser --- les wedges se
réordonnent avec des signes, $\dd$ différencie les
coefficients en de nouveaux $\dd x_j$, les \omterm{def:b3:forms:pullback}{images
réciproques} substituent~; faire confiance à l'algèbre, elle
encode tout jacobien. (2) Pour intégrer une forme sur une
sous-variété~: paramétrer directement, tirer en arrière,
intégrer le coefficient~; l'\omterm{def:b3:forms:orientation}{orientation} est le seul piège ---
vérifier un repère. (3) Pour prouver une identité
d'intégrale, chercher une forme de Stokes~: l'intégrande
est-elle \omterm{def:b3:forms:closedexact}{exacte}~? le domaine est-il un \omterm{def:b3:forms:boundary}{bord}~? (4) Pour
comparer des intégrales sur deux sous-variétés «~parallèles~»,
appliquer Stokes à la région entre elles (l'argument de
déformation, l'\cref{exo:b3:forms:10}). (5) Une intégrale non
nulle d'une \omterm{def:b3:forms:closedexact}{forme fermée} certifie une obstruction
topologique --- pas de primitive, pas de rétraction, pas
d'extension sans zéro~: c'est ainsi que le problème de
week-end tue les rétractions de la boule.
\end{method}

\section{Exercices}

\begin{exercise}[$\star$]\label{exo:b3:forms:1}
Sur $\R^3$, soit $\omega = x\,\dd y - z\,\dd x$ et $\eta =
\dd x + y\,\dd z$. Calculer $\omega\wedge\eta$, $\dd\omega$,
$\dd\eta$ et $\dd(\omega\wedge\eta)$, et vérifier la règle
de Leibniz graduée sur cet exemple.
\end{exercise}

\begin{exercise}[$\star$]\label{exo:b3:forms:2}
Identifier, sur $\R^3$, les trois incarnations de $\dd$~:
pour $f \in \Omega^0$, $\dd f \leftrightarrow \nabla f$~;
pour la forme de travail $\omega_F = F_1\dd x + F_2\dd y +
F_3\dd z$, $\dd\omega_F \leftrightarrow
\operatorname{curl}F$~; pour la forme de flux $\sigma_F =
F_1\,\dd y\wedge\dd z + F_2\,\dd z\wedge\dd x + F_3\,\dd
x\wedge\dd y$, $\dd\sigma_F \leftrightarrow
\operatorname{div}F$. Déduire de $\dd^2 = 0$ les identités
$\operatorname{curl}\nabla f = 0$ et
$\operatorname{div}\operatorname{curl}F = 0$.
\end{exercise}

\begin{exercise}[$\star\star$]\label{exo:b3:forms:3}
Décider si chaque $1$-forme est \omterm{def:b3:forms:closedexact}{fermée}, \omterm{def:b3:forms:closedexact}{exacte} sur son
domaine, et calculer une primitive lorsqu'il en existe une~:
(a) $(2xy + z^2)\,\dd x + x^2\,\dd y + 2xz\,\dd z$ sur
$\R^3$~;
(b) $\dfrac{x\,\dd x + y\,\dd y}{x^2 + y^2}$ sur
$\R^2\setminus\{0\}$~;
(c) $\dfrac{-y\,\dd x + x\,\dd y}{x^2 + y^2}$ sur le
demi-plan $\{x > 0\}$.
\end{exercise}

\begin{exercise}[$\star\star$]\label{exo:b3:forms:4}
(La forme angulaire) Vérifier que $\omega_\theta$
(l'\cref{ex:b3:forms:angular}) est \omterm{def:b3:forms:closedexact}{fermée}~; calculer
$\int_\gamma\omega_\theta$ pour $\gamma$ le cercle de rayon
$r$ autour de $0$~; conclure que $\omega_\theta$ n'est pas
\omterm{def:b3:forms:closedexact}{exacte} sur $\R^2\setminus\{0\}$, et que
$\R^2\setminus\{0\}$ n'est étoilé par rapport à aucun de ses
points (deux voies~: via le~\cref{thm:b3:forms:poincare}, et
directement par la géométrie).
\end{exercise}

\begin{exercise}[$\star\star$]\label{exo:b3:forms:5}
(Forme d'aire d'une hypersurface) Soit $M \subseteq \R^n$
une hypersurface orientée compacte dont l'\omterm{def:b3:forms:orientation}{orientation} est
donnée par un champ de \omterm{def:b3:groups:normal}{normales} unitaires $\nu$ ($(w_1,
\dots, w_{n-1})$ positive ssi $(\nu, w_1, \dots,
w_{n-1})$ positive dans $\R^n$). Montrer que la
$(n-1)$-forme $\sigma_\nu =
\sum_i(-1)^{i-1}\nu_i\,\dd x_1\wedge\dots\wedge\widehat{\dd
x_i}\wedge\dots\wedge\dd x_n$ se restreint sur $M$ à la
\emph{forme d'aire}~: pour une paramétrisation directe
$\gamma$, $\gamma^*\sigma_\nu = \sqrt{\det G}\,\dd
u_1\wedge\dots \wedge\dd u_{n-1}$, où $G = ({}^t
D\gamma)(D\gamma)$ est la matrice de Gram. \emph{(Noter que
$(\gamma^*\sigma_\nu)(e_1, \dots, e_{n-1}) = \det(\nu,
\partial_1\gamma, \dots, \partial_{n-1}\gamma)$ et carrer ce
déterminant.)} Calculer $\int_{S^2}\sigma_\nu = 4\pi$.
\end{exercise}

\begin{exercise}[$\star\star$]\label{exo:b3:forms:6}
Calculer $\int_{S^2}\omega$ pour $\omega = x\,\dd y\wedge\dd
z + y\,\dd z\wedge\dd x + z\,\dd x\wedge\dd y$~: (a)
directement en coordonnées sphériques~; (b) via Stokes sur la
boule unité. En déduire $\operatorname{vol}(B^3) =
\frac{4\pi}3$ à partir de l'aire de $S^2$, et généraliser~:
$n\operatorname{vol}(B^n) = \operatorname{area}(S^{n-1})$,
cohérent avec le~\cref{thm:b3:product:ballvolume}.
\end{exercise}

\begin{exercise}[$\star\star$]\label{exo:b3:forms:7}
(Les dictionnaires) Dériver soigneusement
du~\cref{thm:b3:forms:stokes}~: (a) le théorème de la
divergence dans $\R^3$ (combiner l'\cref{exo:b3:forms:2}
et l'\cref{exo:b3:forms:5})~; (b) le théorème de
Kelvin--Stokes $\int_S\langle\operatorname{curl}F,
\nu\rangle\,\dd S = \oint_{\partial S}\langle F,
\tau\rangle\,\dd\ell$ pour une surface orientée compacte à
\omterm{def:b3:forms:boundary}{bord} dans $\R^3$. Vérifier que les conventions d'\omterm{def:b3:forms:orientation}{orientation}
coïncident sur la demi-sphère supérieure bordée par
l'équateur.
\end{exercise}

\begin{exercise}[$\star\star$]\label{exo:b3:forms:8}
(Identités de Green) Pour $u, v$ lisses sur un \omterm{def:b3:topology:topology}{voisinage} d'un
domaine compact $\Omega \subseteq \R^n$ à \omterm{def:b3:forms:boundary}{bord} lisse,
prouver
\[
\int_\Omega\bigl(u\,\Delta v + \langle\nabla u, \nabla
v\rangle\bigr) = \int_{\partial\Omega}u\,\partial_\nu
v\,\dd S,
\qquad
\int_\Omega\bigl(u\,\Delta v - v\,\Delta u\bigr) =
\int_{\partial\Omega}\bigl(u\,\partial_\nu v -
v\,\partial_\nu u\bigr)\dd S .
\]
En déduire~: une \omterm{prop:b3:conformal:harmonicholo}{fonction harmonique} sur $\Omega$ s'annulant
sur $\partial\Omega$ s'annule identiquement, et deux
\omterm{prop:b3:conformal:harmonicholo}{fonctions harmoniques} aux mêmes valeurs au \omterm{def:b3:forms:boundary}{bord} coïncident
--- unicité dans le problème de Dirichlet
du~\cref{ch:b3:conformal}.
\end{exercise}

\begin{exercise}[$\star\star$]\label{exo:b3:forms:9}
(Changement de variables, forme orientée) Soit
$\varphi\colon U \to V$ un difféomorphisme d'\omterm{def:b3:topology:topology}{ouverts} de
$\R^n$ avec $\det D\varphi > 0$, et $f$ \omterm{def:b3:topology:continuity}{continue} à support
compact dans $V$. Montrer que l'identité d'\omterm{def:b3:forms:pullback}{image réciproque}
$\int_U\varphi^*(f\,\dd x_1\wedge\dots\wedge\dd x_n) =
\int_Vf\,\dd x_1\wedge\dots\wedge\dd x_n$ est équivalente au
théorème de changement de variables
(le~\cref{thm:b3:product:changeofvar}) pour de tels $\varphi$,
et expliquer exactement où est passée la valeur absolue sur
le jacobien.
\end{exercise}

\begin{exercise}[$\star\star\star$]\label{exo:b3:forms:10}
(Déformation) Soit $\omega$ une $2$-forme \emph{\omterm{def:b3:forms:closedexact}{fermée}} sur
$\R^3\setminus\{0\}$ et $S_r$ la sphère de rayon $r$ centrée
en $0$. Montrer que $\int_{S_r}\omega$ ne dépend pas de $r >
0$ \emph{(appliquer Stokes à la coquille entre deux rayons~;
veiller aux deux \omterm{def:b3:forms:orientation}{orientations} de \omterm{def:b3:forms:boundary}{bord})}. Appliquer à la
\emph{forme d'angle solide}
\[
\omega = \frac{x\,\dd y\wedge\dd z + y\,\dd z\wedge\dd x +
z\,\dd x\wedge\dd y}{(x^2 + y^2 + z^2)^{3/2}} :
\]
vérifier qu'elle est \omterm{def:b3:forms:closedexact}{fermée}, calculer $\int_{S_r}\omega =
4\pi$, et conclure qu'elle est \omterm{def:b3:forms:closedexact}{fermée} mais non \omterm{def:b3:forms:closedexact}{exacte} sur
$\R^3\setminus\{0\}$ --- la sœur bidimensionnelle de
$\omega_\theta$, et le \omterm{def:b3:rings:content}{contenu} géométrique de la loi de
Gauss en électrostatique.
\end{exercise}

\begin{exercise}[$\star\star$]\label{exo:b3:forms:11}
Soit $\gamma$ une courbe lisse \omterm{def:b3:forms:closedexact}{fermée} dans
$\R^2\setminus\{0\}$ avec $n = \operatorname{Ind}_\gamma(0)$.
Montrer $\int_\gamma\omega = n\int_{S^1}\omega$ pour
\emph{toute} $1$-forme \omterm{def:b3:forms:closedexact}{fermée} $\omega$ sur
$\R^2\setminus\{0\}$ \emph{(écrire $\omega = c\,\omega_\theta
+ \dd f$ par l'\cref{exo:b3:forms:12})}. Interprétation~: sur le
plan épointé, le \omterm{def:b3:forms:winding}{nombre d'enroulement} est la seule
obstruction à l'annulation des périodes.
\end{exercise}

\begin{exercise}[$\star\star\star$]\label{exo:b3:forms:12}
(Premier calcul de de Rham) Montrer que toute $1$-forme
\omterm{def:b3:forms:closedexact}{fermée} $\omega$ sur $U = \R^2\setminus\{0\}$ s'écrit de
façon unique
\[
\omega = c\,\omega_\theta + \dd f,
\qquad c = \frac1{2\pi}\int_{S^1}\omega,\quad f \in
\mathcal C^\infty(U) :
\]
définir $f(p)$ en intégrant $\omega - c\,\omega_\theta$ le
long d'un chemin de $(1,0)$ à $p$ (morceau radial puis arc
circulaire), montrer que le résultat est indépendant des
choix précisément parce que la période sur $S^1$ s'annule, et
vérifier $\dd f = \omega - c\,\omega_\theta$. Conclure~:
$H^1(\R^2\setminus\{0\}) \cong \R$, engendré par la forme
angulaire.
\end{exercise}

\section{Problème~: le théorème du point fixe de Brouwer}

\begin{problem}[{Problème de week-end --- pas de rétraction,
pas d'évasion}]\label{pb:b3:forms:1}
Le théorème de Brouwer affirme que \emph{toute \omterm{def:b3:topology:continuity}{application
continue} de la boule unité \omterm{def:b3:forms:closedexact}{fermée} $\bar B = \bar B^n
\subseteq \R^n$ dans elle-même a un point fixe} --- l'un des
grands théorèmes des mathématiques, avec des conséquences de
la théorie des jeux (équilibres de Nash) à l'analyse
matricielle. La preuve par \omterm{def:b3:forms:diffform}{formes différentielles} est la plus
propre connue~: Stokes montre que la sphère n'est pas une
rétractée de la boule, et tout s'ensuit. Tout au long, $S =
S^{n-1} = \partial\bar B$, $n \geq 2$, et
\[
\sigma = \sum_{i=1}^n(-1)^{i-1}x_i\,
\dd x_1\wedge\dots\wedge\widehat{\dd x_i}\wedge\dots\wedge
\dd x_n
\]
(le chapeau supprime le facteur).

\medskip
\noindent\textbf{Partie I --- L'instrument de \omterm{def:b3:measure:measure}{mesure}.}
\begin{enumerate}
  \item Calculer $\dd\sigma$, et déduire de Stokes
        (le~\cref{thm:b3:forms:stokes}) que $\int_S\sigma =
        n\operatorname{vol}(\bar B) > 0$, où $S$ porte
        l'\omterm{def:b3:forms:orientation}{orientation} de \omterm{def:b3:forms:boundary}{bord} de la boule.
  \item Pour $n = 2$ et $n = 3$, identifier la restriction de
        $\sigma$ à $S$ avec les formes de longueur d'arc et
        d'aire (l'\cref{exo:b3:forms:5} avec $\nu(x) = x$) et
        recalculer $\int_S\sigma$ directement.
  \item Soit $W \subseteq \R^N$ \omterm{def:b3:topology:topology}{ouvert} et $\varphi\colon W
        \to \R^n$ lisse avec $\norm{\varphi(x)} = 1$ pour
        tout $x \in W$. Montrer que
        $\varphi^*(\dd\sigma) = 0$. \emph{(Différentier
        $\norm\varphi^2 = 1$~: l'image de $D\varphi(x)$
        vit dans l'hyperplan $\varphi(x)^\perp$, de
        dimension $n - 1$~; puis appliquer
        la~\cref{prop:b3:forms:linearpullback}(2).)}
  \item Où est la faille dans la «~preuve~» suivante que
        $\int_S\sigma = 0$~: «~$S$ est compacte sans \omterm{def:b3:forms:boundary}{bord},
        et $\sigma$ restreinte à $S$ est une forme de \omterm{def:b3:galois:extension}{degré}
        \omterm{def:b3:rings:primemaximal}{maximal} dessus, donc \omterm{def:b3:forms:closedexact}{fermée}, donc $\int_S\sigma =
        \int_S\dd(\text{quelque chose}) = 0$ par Stokes~»~?
        (Pointer le mot faux.)
  \item Expliquer en un paragraphe la stratégie de la
        Partie~II~: que va-t-on intégrer, sur quoi, et d'où
        viendra la contradiction.
\end{enumerate}

\medskip
\noindent\textbf{Partie II --- Pas de rétraction lisse.}
Supposer, par l'absurde, que $r$ est une \emph{rétraction
lisse} de la boule sur sa sphère~: $r$ est lisse sur un
\omterm{def:b3:topology:topology}{voisinage} de $\bar B$, $r(\bar B) \subseteq S$, et $r(x) =
x$ pour tout $x \in S$.
\begin{enumerate}[resume]
  \item Justifier $\int_Sr^*\sigma = \int_S\sigma$
        \emph{(sur $S$, $r$ se restreint à l'identité~: si
        $\gamma$ est une paramétrisation directe d'un
        morceau de $S$, alors $r\circ\gamma = \gamma$)}.
  \item En utilisant Stokes sur $\bar B$, montrer
        $\int_Sr^*\sigma = \int_{\bar B}\dd(r^*\sigma)$.
  \item Montrer $\dd(r^*\sigma) = r^*(\dd\sigma) = 0$
        (question~3 appliquée à $\varphi = r$), et
        conclure~: \emph{il n'existe pas de rétraction lisse
        $\bar B \to S$}.
  \item Régler le cas exclu $n = 1$ à la main~: montrer
        directement qu'aucune \omterm{def:b3:topology:continuity}{application continue}
        $\intcc{-1}1 \to \{-1, 1\}$ ne fixe les deux
        extrémités, et nommer le théorème utilisé.
\end{enumerate}

\medskip
\noindent\textbf{Partie III --- Brouwer lisse.} Soit $g$
lisse sur un \omterm{def:b3:topology:topology}{voisinage} de $\bar B$ avec $g(\bar B)
\subseteq \bar B$ et \emph{aucun} point fixe dans $\bar B$.
\begin{enumerate}[resume]
  \item Montrer $\delta = \min_{x\in\bar B}\norm{g(x) - x} >
        0$.
  \item Pour $x \in \bar B$ soit $u(x) = \frac{x -
        g(x)}{\norm{x - g(x)}}$ et soit $r(x) = x +
        t(x)\,u(x)$ l'intersection du rayon $\{x + tu(x) : t
        \geq 0\}$ avec $S$. Résoudre le quadratique
        $\norm{x + tu}^2 = 1$ et obtenir
        \[
        t(x) = -\langle x, u(x)\rangle +
        \sqrt{1 - \norm x^2 + \langle x, u(x)\rangle^2}
        \;\geq\; 0 .
        \]
  \item Montrer que la quantité sous le radical est \emph{strictement}
        positif sur $\bar B$~: si $1 - \norm x^2 + \langle
        x, u(x)\rangle^2 = 0$ alors $\norm x = 1$ et
        $\langle x, u(x)\rangle = 0$, i.e.\ $\langle x, x -
        g(x)\rangle = 0$, i.e.\ $\langle x, g(x)\rangle =
        1$~; par Cauchy--Schwarz avec $\norm x = 1$,
        $\norm{g(x)} \leq 1$, cela force $g(x) = x$ ---
        exclu. En déduire que $r$ est lisse sur un \omterm{def:b3:topology:topology}{voisinage}
        de $\bar B$.
  \item Montrer $r(\bar B) \subseteq S$ et $r(x) = x$ pour
        $x \in S$ \emph{(pour $\norm x = 1$, vérifier $t(x)
        = 0$ en utilisant $\langle x, u(x)\rangle \geq 0$,
        qui suit de $\langle x, x - g(x)\rangle = 1 -
        \langle x, g(x)\rangle \geq 0$)}. Conclure avec la
        Partie~II~: \emph{toute application lisse de $\bar
        B$ dans elle-même a un point fixe.}
\end{enumerate}

\medskip
\noindent\textbf{Partie IV --- Brouwer continu.} Soit
$f\colon\bar B\to\bar B$ \omterm{def:b3:topology:continuity}{continue} sans point fixe.
\begin{enumerate}[resume]
  \item Montrer $\varepsilon = \min_{\bar B}\norm{f - x} >
        0$, et produire une application
        \emph{polynomiale} $p\colon\R^n\to \R^n$ avec
        $\sup_{\bar B}\norm{p - f} < \varepsilon/2$
        (Stone--Weierstrass,
        le~\cref{thm:b3:complete:stoneweierstrass}, coordonnée
        par coordonnée --- justifier le passage de
        l'approximation scalaire à vectorielle).
  \item L'application $p$ peut quitter la boule~; poser $g =
        \frac{p}{1 + \varepsilon/2}$. Montrer $g(\bar B)
        \subseteq \bar B$ et $\sup_{\bar B}\norm{g - f} <
        \varepsilon$.
  \item En déduire une contradiction avec la Partie~III et
        conclure~: \emph{toute \omterm{def:b3:topology:continuity}{application continue} $\bar
        B^n \to \bar B^n$ a un point fixe.}
  \item Montrer par des exemples que le théorème échoue sur~:
        la boule ouverte~; la sphère $S$~; un anneau fermé.
        Quelle propriété de $\bar B$ chaque contre-exemple
        perd-il~?
\end{enumerate}

\medskip
\noindent\textbf{Partie V --- Dividendes.}
\begin{enumerate}[resume]
  \item (Perron--Frobenius, existence) Soit $A$ une matrice
        $n\times n$ à entrées toutes $> 0$, et $\Delta =
        \{x \in \R^n : x_i \geq 0,\ \sum x_i = 1\}$.
        Montrer que l'application $x \mapsto
        Ax/\norm{Ax}_1$ est bien définie et \omterm{def:b3:topology:continuity}{continue} sur
        $\Delta$, que $\Delta$ est homéomorphe à une boule
        \omterm{def:b3:forms:closedexact}{fermée} de $\R^{n-1}$ (\omterm{def:b3:topology:continuity}{homéomorphisme} radial depuis un
        compact convexe d'intérieur non vide dans son
        enveloppe affine), et conclure que $A$ a un vecteur
        propre à entrées strictement positives et valeur
        propre $> 0$.
  \item En déduire que toute matrice stochastique à entrées
        positives (colonnes de somme $1$) a un vecteur de
        probabilité stationnaire $\pi = A\pi$ --- le
        vecteur de type PageRank. (L'unicité vaut aussi mais
        demande d'autres outils.)
  \item (Boule chevelue, mise en place) Soit $v$ lisse sur un
        \omterm{def:b3:topology:topology}{voisinage} de $S = S^{n-1}$ avec $\langle v(x),
        x\rangle = 0$ et $\norm{v(x)} = 1$ pour $x \in S$
        (un \emph{champ tangent unitaire}). Pour $t \in \R$
        poser $F_t(x) = x + t\,v(x)$. Montrer
        $\norm{F_t(x)} = \sqrt{1 + t^2}$ sur $S$~: $F_t$
        envoie $S$ dans la sphère $\sqrt{1+t^2}\,S$.
  \item Montrer que $P(t) = \int_SF_t^*\sigma$ est un
        \emph{polynôme} en $t$ \emph{(chaque fonction
        coefficient de $F_t^*\sigma$ dans une carte est
        polynomiale en $t$, à coefficients lisses dans la
        variable de carte~; l'intégration est linéaire)}.
  \item Montrer que pour $\abs t$ petit, $F_t$ est un
        difféomorphisme de $S$ sur $\sqrt{1+t^2}\,S$~:
        injectivité pour $t\operatorname{Lip}(v) < 1$~;
        difféomorphisme local par le théorème d'inversion
        locale (le~\cref{thm:b3:submanifolds:ift} en cartes)~;
        image ouverte et \omterm{def:b3:forms:closedexact}{fermée} dans la sphère cible
        \omterm{def:b3:topology:connected}{connexe}. En déduire, en utilisant
        le~\cref{lem:b3:forms:consistency} et le
        redimensionnement $\sigma_{\lambda x} =
        \lambda^{n}\,\sigma_x$ sous $x \mapsto \lambda x$
        (le vérifier), que
        \[
        P(t) = \pm(1 + t^2)^{n/2}\int_S\sigma,
        \qquad\text{avec le signe } + \text{ pour } t
        \text{ petit}
        \]
        (\omterm{def:b3:forms:orientation}{orientation} préservée par continuité depuis $t =
        0$).
  \item Conclure (Milnor)~: si $n$ est \emph{impair}, $(1 +
        t^2)^{n/2}$ n'est pas un polynôme en $t$, or il
        coïncide avec le polynôme $P(t)/\int_S\sigma$ près
        de $0$ --- contradiction. D'où \emph{les sphères de
        dimension paire $S^{n-1}$ ($n$ impair) ne portent
        aucun champ tangent unitaire}, et, en normalisant et
        lissant (convoluer composante par composante et
        projeter --- justifier les deux étapes), aucun champ
        tangent continu nulle part nul~: tout vent sur Terre
        laisse un point calme.

  \item (Les sphères impaires se peignent librement)
        Exhiber sur $S^{2m-1} \subseteq \R^{2m} \cong \C^m$
        un champ tangent unitaire lisse explicite~: $v(x) =
        \iu x$ en notation complexe, i.e.
        \[
        v(x_1, y_1, \dots, x_m, y_m) =
        (-y_1, x_1, \dots, -y_m, x_m) .
        \]
        Vérifier la tangence et la longueur unité, et
        conclure que la dichotomie de parité de la
        question~23 est tranchante~: une sphère est peignable
        exactement lorsque sa dimension est impaire. Où
        l'argument polynomial de la question~22 casse-t-il
        pour $n$ pair~?
  \item (Zéros issus du comportement au \omterm{def:b3:forms:boundary}{bord}) Soit $f \colon
        \bar B^n \to \R^n$ \omterm{def:b3:topology:continuity}{continue} avec $\langle f(x),
        x\rangle \geq 0$ pour tout $x \in S^{n-1}$. Montrer
        que $f$ s'annule quelque part dans $\bar B^n$.
        \emph{(Sinon, $g(x) = -f(x)/\norm{f(x)}$ envoie
        $\bar B$ continûment dans $S \subseteq \bar B$~;
        appliquer Brouwer à $g$ et contredire l'hypothèse
        au \omterm{def:b3:forms:boundary}{bord}.)} En déduire le \emph{critère de
        surjectivité}~: une \omterm{def:b3:topology:continuity}{application continue}
        $F\colon\R^n\to\R^n$ avec $\frac{\langle F(x),
        x\rangle}{\norm x} \to +\infty$ lorsque $\norm x
        \to \infty$ est surjective --- l'ancêtre de
        dimension finie des arguments de coercivité de
        l'analyse non linéaire.
\end{enumerate}
\end{problem}
